\documentclass[preprint]{revtex4-2}

\usepackage{amsmath,amssymb,amsthm}
\usepackage{graphicx}
\usepackage{hyperref}
\usepackage{booktabs}

\newtheorem{theorem}{Theorem}
\newtheorem{corollary}{Corollary}

\newcommand{\rad}{\operatorname{rad}}
\newcommand{\sopfr}{\operatorname{sopfr}}
\DeclareMathOperator{\lcm}{lcm}

\begin{document}

\title{Consonance, Crystals, and Orbits:\\
The $\varphi(n) \leq 2$ Filter Across Domains}

\author{Park, Min Woo}
\affiliation{Independent Researcher}
\email{nerve011235@gmail.com}

\date{\today}

\begin{abstract}
Why can crystals display only 1-, 2-, 3-, 4-, and 6-fold rotational symmetry? Why do the perfect consonances of Western music use only the frequency ratios $1{:}1$, $2{:}1$, $3{:}2$, and $4{:}3$? Why does the chromatic scale have 12 notes? These three questions, drawn from solid-state physics, psychoacoustics, and music theory, appear to have nothing in common. We show that they share a single number-theoretic answer: the Euler totient condition $\varphi(n) \leq 2$, which selects precisely the set $\{1, 2, 3, 4, 6\}$. The crystallographic restriction demands that $2\cos(2\pi/n)$ be an integer, which is equivalent to the cyclotomic polynomial $\Phi_n(x)$ having degree at most~2, i.e., $\varphi(n) \leq 2$. Musical consonance, in the Helmholtz--Euler tradition, privileges small-integer frequency ratios, and the ratios built from the divisors of~6 are exactly those with numerator and denominator in $\{1,2,3,4,6\}$. The number $12 = \lcm(2,3,4,6)$ then emerges as the smallest period accommodating all consonant intervals. We further observe that the perfect-number identity $\sigma_{-1}(6) = 2$ encodes the factorization of the octave as fifth~$\times$~fourth: $(3/2)(4/3) = 2$. A weaker but suggestive connection links the same set to the innermost stable circular orbit (ISCO) in Schwarzschild spacetime at $r = 6M$, where $6 = 2 \times 3$ and $L^2_{\text{ISCO}} = 12M^2$. The unifying algebraic object is the prime pair~$(2,3)$, uniquely characterized by $(p-1)(q-1) = 2$, which forces $\varphi(pq) = 2$ and makes $pq = 6$ the largest member of the $\varphi(n) \leq 2$ family. All proofs are elementary and self-contained.
\end{abstract}

\keywords{Euler totient, crystallographic restriction, musical consonance, perfect numbers, prime pair, Schwarzschild ISCO}

\pacs{02.10.De, 61.50.Ah, 43.75.Bc, 04.20.Cv}

\maketitle

\section{Introduction}

Consider three questions from three different fields:

\begin{enumerate}
\item \textbf{Crystallography.} Why do two-dimensional lattices admit only rotational symmetries of order 1, 2, 3, 4, and 6 --- but never 5, 7, 8, or any higher order?

\item \textbf{Music theory.} Why are the ``perfect'' consonances --- unison, octave, fifth, and fourth --- built from the ratios $1{:}1$, $2{:}1$, $3{:}2$, and $4{:}3$, using only the numbers 1 through 4 (and, by extension, 6)?

\item \textbf{Scale construction.} Why does virtually every tuning tradition converge on a division of the octave into 12 semitones?
\end{enumerate}

At first glance these appear to be entirely separate phenomena. The crystallographic restriction is a consequence of lattice geometry. Consonance is grounded in the physiology of the auditory system. And the 12-note chromatic scale is, one might assume, a historical accident of Western music.

We will show that all three phenomena are controlled by the same number-theoretic filter: the condition $\varphi(n) \leq 2$ on the Euler totient function, which admits precisely the integers $n \in \{1, 2, 3, 4, 6\}$. The argument requires nothing beyond elementary number theory and linear algebra.

We also discuss a fourth, more speculative connection: the innermost stable circular orbit (ISCO) in Schwarzschild general relativity occurs at $r = 6M$, and the critical angular momentum satisfies $L^2 = 12M^2 = \lcm(2,3,4,6) \cdot M^2$. While the ISCO derivation does not directly invoke the totient condition, the appearance of the same numbers is at least noteworthy.

\textbf{Outline.} Section~\ref{sec:theorem} proves the totient filter theorem. Sections~\ref{sec:crystal}, \ref{sec:music}, and~\ref{sec:orbits} apply it to crystallography, music, and orbital mechanics respectively. Section~\ref{sec:prime-pair} identifies the algebraic root: the prime pair~$(2,3)$. Section~\ref{sec:discussion} discusses scope and limitations.

\section{The $\varphi(n) \leq 2$ Theorem}\label{sec:theorem}

\subsection{Statement and proof}

\begin{theorem}\label{thm:totient}
The Euler totient function satisfies $\varphi(n) \leq 2$ if and only if $n \in \{1, 2, 3, 4, 6\}$.
\end{theorem}

\begin{proof}
Direct computation gives:

\begin{center}
\begin{tabular}{l*{11}{c}}
\toprule
$n$ & 1 & 2 & 3 & 4 & 5 & 6 & 7 & 8 & 9 & 10 & 12 \\
$\varphi(n)$ & 1 & 1 & 2 & 2 & 4 & 2 & 6 & 4 & 6 & 4 & 4 \\
\bottomrule
\end{tabular}
\end{center}

So the ``if'' direction is immediate. For the ``only if,'' we must show $\varphi(n) \geq 3$ for all $n \geq 7$.

Write $n = p_1^{a_1} \cdots p_k^{a_k}$ with $p_1 < \cdots < p_k$. Then
\begin{equation}
\varphi(n) = \prod_{i=1}^{k} p_i^{a_i - 1}(p_i - 1).
\end{equation}

\textbf{Case~1: $n = p^a$ is a prime power.}
Then $\varphi(n) = p^{a-1}(p-1)$. For $p \geq 5$: $\varphi(p) = p - 1 \geq 4 > 2$. For $p = 3$, $a \geq 2$: $\varphi(9) = 6 > 2$. For $p = 2$, $a \geq 3$: $\varphi(8) = 4 > 2$. The only prime powers with $\varphi \leq 2$ are $2$, $4$, $3$.

\textbf{Case~2: $n$ has at least two distinct prime factors.}
Then $\varphi(n) \geq (p_1 - 1)(p_2 - 1)$.
\begin{itemize}
\item If $p_1 = 2, p_2 = 3$: $(p_1 - 1)(p_2 - 1) = 2$, so $\varphi(n) \geq 2$. But if any $a_i \geq 2$ or there is a third prime factor $p_3 \geq 5$, additional factors push $\varphi(n) \geq 4$. So $\varphi(n) = 2$ requires $n = 6$.
\item If $p_1 = 2, p_2 \geq 5$: $(p_1 - 1)(p_2 - 1) \geq 4 > 2$.
\item If $p_1 \geq 3, p_2 \geq 5$: $(p_1 - 1)(p_2 - 1) \geq 8 > 2$.
\end{itemize}

Combining all cases, $\varphi(n) \leq 2$ if and only if $n \in \{1, 2, 3, 4, 6\}$.
\end{proof}

\subsection{Connection to cyclotomic polynomials}

The $n$-th cyclotomic polynomial $\Phi_n(x)$ is the minimal polynomial of $e^{2\pi i/n}$ over $\mathbb{Q}$, and $\deg \Phi_n = \varphi(n)$. Therefore:
\begin{equation}
\varphi(n) \leq 2 \quad \Longleftrightarrow \quad \Phi_n(x) \text{ is linear or quadratic over } \mathbb{Q}.
\end{equation}

Explicitly:

\begin{center}
\begin{tabular}{ccc}
\toprule
$n$ & $\Phi_n(x)$ & $\deg$ \\
\midrule
1 & $x - 1$ & 1 \\
2 & $x + 1$ & 1 \\
3 & $x^2 + x + 1$ & 2 \\
4 & $x^2 + 1$ & 2 \\
6 & $x^2 - x + 1$ & 2 \\
\bottomrule
\end{tabular}
\end{center}

For $n \geq 7$, $\Phi_n$ has degree~$\geq 4$, and in particular $\cos(2\pi/n)$ is not a rational number. This is the algebraic fact underlying the crystallographic restriction.

\section{Application 1: Crystallographic Restriction}\label{sec:crystal}

\subsection{The constraint}

A rotation by angle $\theta = 2\pi/n$ is compatible with a two-dimensional lattice if and only if the rotation matrix
\begin{equation}
R = \begin{pmatrix} \cos\theta & -\sin\theta \\ \sin\theta & \cos\theta \end{pmatrix}
\end{equation}
maps lattice vectors to lattice vectors. Since lattice vectors form a $\mathbb{Z}$-module of rank~2, the matrix~$R$ must have integer trace:
\begin{equation}
\mathrm{tr}(R) = 2\cos(2\pi/n) \in \mathbb{Z}.
\end{equation}

\subsection{Why this is $\varphi(n) \leq 2$}

The number $2\cos(2\pi/n)$ is a root of the minimal polynomial of $\zeta_n + \zeta_n^{-1}$ over $\mathbb{Q}$, whose degree is $\varphi(n)/2$ for $n \geq 3$. For $2\cos(2\pi/n)$ to be a rational integer, this degree must be~1, requiring $\varphi(n) \leq 2$.

By Theorem~\ref{thm:totient}, the allowed orders are $n \in \{1, 2, 3, 4, 6\}$, giving:

\begin{center}
\begin{tabular}{ccccc}
\toprule
$n$ & 1 & 2 & 3 & 4 & 6 \\
\midrule
$\theta$ (degrees) & 360 & 180 & 120 & 90 & 60 \\
$\mathrm{tr}(R)$ & 2 & $-2$ & $-1$ & 0 & 1 \\
\bottomrule
\end{tabular}
\end{center}

This is the \textbf{crystallographic restriction theorem}. It explains why snowflakes are hexagonal ($n=6$), why square and hexagonal tilings exist but pentagonal ones do not, and why the 32 crystallographic point groups in 3D are built from these rotational orders.

\subsection{Figure: the five allowed symmetries}

\begin{figure}[h]
\centering
% Placeholder for figure of the five allowed symmetry types
\fbox{\parbox{0.8\columnwidth}{\centering\vspace{2em}
$n=1$: identity \quad $n=2$: inversion \quad $n=3$: triangle\\[1em]
$n=4$: square \quad $n=6$: hexagon\\[1em]
\textit{(Five allowed lattice symmetry orders)}
\vspace{1em}}}
\caption{The five rotational symmetry orders $n \in \{1,2,3,4,6\}$ compatible with lattice periodicity. The five-fold case $n=5$ fails because $2\cos(72^\circ) = (\sqrt{5}-1)/2 \approx 0.618$ is irrational.}
\label{fig:symmetries}
\end{figure}

\section{Application 2: Musical Consonance}\label{sec:music}

\subsection{The Helmholtz--Euler theory}

Helmholtz (1863) and Euler (before him) recognized that the ear perceives two tones as consonant when their frequency ratio is a ratio of small integers~\cite{helmholtz}. The ranking of intervals by consonance is approximately:

\begin{center}
\begin{tabular}{lcc}
\toprule
Interval & Ratio & Consonance rank \\
\midrule
Unison & $1{:}1$ & 1 (most consonant) \\
Octave & $2{:}1$ & 2 \\
Perfect fifth & $3{:}2$ & 3 \\
Perfect fourth & $4{:}3$ & 4 \\
Major sixth & $5{:}3$ & 5 \\
Major third & $5{:}4$ & 6 \\
Minor third & $6{:}5$ & 7 \\
\bottomrule
\end{tabular}
\end{center}

The top four --- the \textbf{perfect consonances} recognized since antiquity --- use only the integers $\{1, 2, 3, 4\}$ in their ratios, all of which are divisors of~6. The remaining ``imperfect'' consonances require the prime~5.

\subsection{The divisors of 6 and frequency ratios}

The \textbf{distinct reduced ratios} within one octave from the set $\{1,2,3,4,6\}$ are precisely: $2/1$, $3/2$, $4/3$ --- the three nontrivial perfect consonances. The set $\{1,2,3,4,6\}$ generates exactly the intervals that every musical tradition treats as most consonant.

\subsection{The octave identity and $\sigma_{-1}(6) = 2$}

The sum of reciprocals of the divisors of~6 is:
\begin{equation}
\sigma_{-1}(6) = \frac{1}{1} + \frac{1}{2} + \frac{1}{3} + \frac{1}{6} = 2.
\end{equation}
This is the defining property of a perfect number: $\sigma_{-1}(n) = 2$ if and only if $\sigma(n) = 2n$.

The product of the two nontrivial consonances within one octave is:
\begin{equation}
\frac{3}{2} \times \frac{4}{3} = \frac{12}{6} = 2 = \text{octave}.
\end{equation}
The \textbf{fifth composed with the fourth equals the octave}. The identity $1/2 + 1/3 + 1/6 = 1$ (the sum of proper reciprocal divisors) means that the ``cost'' of each consonance, measured as a fraction of the octave, adds up exactly to the whole. This additive closure is unique to~6 among all positive integers.

\subsection{Why 12 notes?}

The chromatic scale divides the octave into 12 equal semitones. Algebraically, we want $N$ such that $2^{k/N} \approx 3/2$ for some integer~$k$. The convergents of $\log_2(3/2) \approx 0.58496$ are $1/2, 3/5, 7/12, 24/41, \ldots$ The convergent $7/12$ gives $2^{7/12} = 1.49831 \approx 1.5$, an error of only~0.11\%.

But there is a cleaner number-theoretic explanation:
\begin{equation}
12 = \lcm(2, 3, 4, 6) = \lcm\{n : \varphi(n) \leq 2\} \setminus \{1\}.
\end{equation}
The number~12 is the least common multiple of all nontrivial members of the $\varphi(n) \leq 2$ family --- the smallest period simultaneously compatible with all allowed symmetry orders.

\subsection{Figure: the circle of fifths}

\begin{figure}[h]
\centering
% Placeholder for circle-of-fifths diagram
\fbox{\parbox{0.6\columnwidth}{\centering\vspace{2em}
C -- G -- D -- A -- E -- B\\
$|$ \hspace{6em} $|$\\
F -- B$\flat$ -- E$\flat$ -- A$\flat$ -- D$\flat$ -- G$\flat$/F$\sharp$\\[1em]
12 notes by fifths ($3{:}2$)\\
7 semitones per fifth: $7/12$ of an octave\\
$(3/2)^{12} \approx 2^7$
\vspace{1em}}}
\caption{The 12-note chromatic scale arranged by fifths. The 12-note equal temperament gives the best simultaneous approximation to intervals generated by the primes 2 and 3.}
\label{fig:circle}
\end{figure}

\section{Application 3: Orbital Stability (Weaker Connection)}\label{sec:orbits}

\subsection{The Schwarzschild ISCO}

In general relativity, a test particle orbiting a Schwarzschild black hole of mass~$M$ moves in the effective potential (in units with $G = c = 1$):
\begin{equation}
V_{\text{eff}}(r) = -\frac{M}{r} + \frac{L^2}{2r^2} - \frac{ML^2}{r^3}
\end{equation}
where $L$ is the specific angular momentum.

Circular orbits exist where $V_{\text{eff}}'(r) = 0$:
\begin{equation}
\frac{M}{r^2} - \frac{L^2}{r^3} + \frac{3ML^2}{r^4} = 0 \implies L^2 = \frac{Mr^2}{r - 3M}.
\end{equation}

The innermost stable circular orbit (ISCO) occurs where $V_{\text{eff}}''(r) = 0$. Substituting and simplifying yields:
\begin{equation}
r(r - 6M) = 0.
\end{equation}

Therefore:
\begin{equation}\label{eq:isco}
\boxed{r_{\text{ISCO}} = 6M, \qquad L^2_{\text{ISCO}} = 12M^2.}
\end{equation}

\subsection{The numbers 6 and 12 again}

The ISCO radius is $6M$ and the critical angular momentum squared is $12M^2$ --- exactly the numbers from the totient filter ($\max\{n : \varphi(n) \leq 2\} = 6$) and the chromatic scale ($\lcm(2,3,4,6) = 12$).

\textbf{We must be honest about the status of this connection.} The ISCO derivation is a specific calculation in Schwarzschild geometry. The number~6 arises from the algebraic structure of a cubic equation, not from a totient condition. What we can say is:

\begin{enumerate}
\item The number~6 appears as a stability boundary in GR, just as it appears as the largest lattice-compatible symmetry order.
\item The number~12 appears as $L^2_{\text{ISCO}}/M^2$, just as it appears as the natural chromatic period.
\item Both 6 and 12 are built exclusively from the primes 2 and 3.
\end{enumerate}

Whether there is a deeper connection remains an open question. We flag this as \textbf{suggestive but not proven} to share the totient origin of the other two applications.

\subsection{Figure: the effective potential}

\begin{figure}[h]
\centering
% Placeholder for effective potential plot
\fbox{\parbox{0.7\columnwidth}{\centering\vspace{2em}
$V_{\text{eff}}(r)$ vs.\ $r/M$\\[1em]
$r = 3M$: photon sphere (unstable null orbits)\\
$r = 6M$: ISCO (marginally stable timelike orbits)\\
$r = 12M$: $L^2_{\text{ISCO}}$ intersects Newtonian prediction
\vspace{1em}}}
\caption{Schematic of the Schwarzschild effective potential showing the ISCO at $r = 6M$ and the photon sphere at $r = 3M$.}
\label{fig:potential}
\end{figure}

\section{The Unifying Object: The Prime Pair $(2, 3)$}\label{sec:prime-pair}

\subsection{The unique pair with $(p-1)(q-1) = 2$}

\begin{theorem}\label{thm:unique-pair}
The primes $p = 2$ and $q = 3$ are the unique pair of distinct primes satisfying $(p-1)(q-1) = 2$.
\end{theorem}

\begin{proof}
Let $p < q$ be distinct primes. Then $p \geq 2$ and $q \geq 3$, so $(p-1)(q-1) \geq 1 \cdot 2 = 2$. Equality holds iff $p-1 = 1$ and $q-1 = 2$, i.e., $p = 2, q = 3$.
\end{proof}

Since $\varphi(pq) = (p-1)(q-1)$ for distinct primes, this immediately gives:

\begin{corollary}
$n = 6 = 2 \cdot 3$ is the unique squarefree semiprime with $\varphi(n) \leq 2$, and consequently the largest element of $\{n : \varphi(n) \leq 2\}$.
\end{corollary}

\subsection{Why $(2, 3)$ controls everything}

The prime pair~$(2, 3)$ sits at the foundation of three structures:

\begin{center}
\begin{tabular}{llll}
\toprule
Structure & Role of 2 & Role of 3 & Combined \\
\midrule
Lattice symmetry & $C_2$ reflection & $C_3$ triangular & $C_6$ hexagonal \\
Musical interval & Octave ($2{:}1$) & Fifth-related ($3{:}2$) & All perfect consonances \\
Orbital mechanics & Binary (attract/repel) & GR cubic correction & ISCO at $r = 6M$ \\
Arithmetic & First prime & Second prime & $\varphi(6) = 2$ (perfect) \\
\bottomrule
\end{tabular}
\end{center}

The pair~$(2,3)$ is special because it is the only consecutive integer pair that is also a prime pair, and $2 \cdot 3 = 6$ is the only number that is both a primorial ($p_2\#$) and a perfect number.

\subsection{The role of $\sigma_{-1}(6) = 2$}

The sum of reciprocal divisors ties the arithmetic to the applications:
\begin{equation}
\sigma_{-1}(6) = \frac{1}{1} + \frac{1}{2} + \frac{1}{3} + \frac{1}{6} = 2.
\end{equation}

Decompose as:
\begin{equation}
\underbrace{\frac{1}{1}}_{\text{unison}} + \underbrace{\frac{1}{2}}_{\text{octave}^{-1}} + \underbrace{\frac{1}{3}}_{\text{twelfth}^{-1}} + \underbrace{\frac{1}{6}}_{\text{product}} = 2 = \text{octave}.
\end{equation}

The identity $1/2 + 1/3 + 1/6 = 1$ (the sum of proper reciprocal divisors) means that the ``cost'' of each consonance adds up exactly to the whole. This additive closure is unique to~6 among all positive integers.

\section{Discussion}\label{sec:discussion}

\subsection{What the totient filter explains}

The condition $\varphi(n) \leq 2$ provides a \textbf{uniform explanation} for:
\begin{enumerate}
\item The crystallographic restriction.
\item The primacy of perfect consonances.
\item The 12-note chromatic scale ($12 = \lcm$ of the filtered set).
\item The octave identity ($\sigma_{-1}(6) = 2$ encodes fifth~$\times$~fourth~$=$~octave).
\end{enumerate}

\subsection{What it does not explain}

The ISCO connection remains \textbf{correlational rather than causal}. The number~6 in $r_{\text{ISCO}} = 6M$ arises from the specific algebraic structure of the Schwarzschild effective potential, not from a totient condition. A genuine unification would require showing that the Schwarzschild stability analysis is an instance of some general ``periodic-system constraint'' that reduces to the totient condition. We are not aware of such a result and do not claim one.

\subsection{Pedagogical value}

The $\varphi(n) \leq 2$ filter is valuable as a teaching tool because it connects three subjects with a single two-line proof, demonstrates how the same algebraic constraint can manifest in physically distinct domains, and is accessible to anyone who knows the definition of the Euler totient function.

\subsection{Relation to prior work}

The crystallographic restriction theorem is classical (Ha\"uy 1822, Frankenheim 1842) and appears in every solid-state physics textbook~\cite{ashcroft,kittel}. Its connection to cyclotomic polynomials is well known in algebra~\cite{lang}. The Helmholtz theory of consonance dates to 1863~\cite{helmholtz}. The observation that $\sigma_{-1}(6) = 2$ encodes the octave decomposition appears to be less widely noted.

What is new here is the explicit identification of $\varphi(n) \leq 2$ as the \textbf{single condition} underlying all three domains, and the observation that $\sigma_{-1}(6) = 2$ gives a number-theoretic expression of the fundamental identity of Western harmony.

\section{Summary}

We have shown that the Euler totient condition $\varphi(n) \leq 2$, which selects the set $\{1, 2, 3, 4, 6\}$, is the common root of the crystallographic restriction theorem, the theory of perfect musical consonance, and the 12-note chromatic scale. The algebraic engine is the prime pair~$(2,3)$, uniquely characterized by $(p-1)(q-1) = 2$, whose product~6 is both the largest totient-filtered integer and the first perfect number. The identity $\sigma_{-1}(6) = (3/2)(4/3) = 2$ equates the perfect-number condition with the octave decomposition into fifth and fourth. A weaker but suggestive parallel links the same numbers to the ISCO in Schwarzschild spacetime. The entire argument is elementary, self-contained, and accessible to undergraduates.

\bigskip
\noindent\textbf{Data availability.} All computational verification scripts are available at \url{https://github.com/need-singularity/TECS-L/tree/main/verify/}.

\begin{acknowledgments}
The author thanks the open-source mathematical community and the OEIS contributors.
\end{acknowledgments}

\appendix

\section{Computational Verification}\label{app:code}

The following Python script verifies the main claims:

\begin{verbatim}
from math import gcd
from functools import reduce

def euler_totient(n):
    return sum(1 for k in range(1, n+1)
               if gcd(k, n) == 1)

def sigma_neg1(n):
    return sum(1/d for d in range(1, n+1)
               if n % d == 0)

def lcm(a, b):
    return a * b // gcd(a, b)

# Verify phi(n) <= 2 filter
for n in range(1, 25):
    phi = euler_totient(n)
    flag = "  <-- PASS" if phi <= 2 else ""
    print(f"{n:2d}: {phi:5d}   {flag}")

# Verify sigma_{-1}(6) = 2
print(f"sigma_{{-1}}(6) = {sigma_neg1(6)}")

# Verify LCM
filtered = [2, 3, 4, 6]
L = reduce(lcm, filtered)
print(f"lcm(2,3,4,6) = {L}")
\end{verbatim}

\section{The 32 Crystallographic Point Groups}\label{app:crystal}

The 32 point groups in three dimensions are generated by rotations of order $n \in \{1, 2, 3, 4, 6\}$ combined with reflections and inversions. They organize into 7 crystal systems:

\begin{center}
\begin{tabular}{lll}
\toprule
Crystal System & Required Symmetry & Allowed $n$ \\
\midrule
Triclinic & $C_1$ or $C_i$ & 1 \\
Monoclinic & $C_2$ axis & 2 \\
Orthorhombic & 3 mutually perp.\ $C_2$ & 2 \\
Tetragonal & $C_4$ axis & 4, 2 \\
Trigonal & $C_3$ axis & 3 \\
Hexagonal & $C_6$ axis & 6, 3, 2 \\
Cubic & 4 $C_3$ axes & 3, 2 \\
\bottomrule
\end{tabular}
\end{center}

Every entry in the ``Allowed~$n$'' column is a member of $\{1, 2, 3, 4, 6\}$. The 230 space groups in 3D (Fedorov 1891, Schoenflies 1891) are obtained by combining these point groups with the 14 Bravais lattices, and no additional rotational orders appear.

\begin{thebibliography}{10}

\bibitem{helmholtz}
H.~von Helmholtz, \emph{On the Sensations of Tone as a Physiological Basis for the Theory of Music} (1863; English trans.\ A.\,J.~Ellis, 1875).

\bibitem{ashcroft}
N.\,W.~Ashcroft and N.\,D.~Mermin, \emph{Solid State Physics} (Harcourt, 1976), Ch.~7.

\bibitem{kittel}
C.~Kittel, \emph{Introduction to Solid State Physics}, 8th ed.\ (Wiley, 2004), Ch.~1.

\bibitem{lang}
S.~Lang, \emph{Algebra}, 3rd ed.\ (Springer, 2002), Ch.~IV.

\bibitem{hardy-wright}
G.\,H.~Hardy and E.\,M.~Wright, \emph{An Introduction to the Theory of Numbers}, 6th ed.\ (Oxford, 2008), Ch.~XVIII.

\bibitem{carroll}
S.~Carroll, \emph{Spacetime and Geometry: An Introduction to General Relativity} (Cambridge, 2019), Ch.~5.

\bibitem{euler}
L.~Euler, ``De numeris amicabilibus,'' \emph{Opera Omnia}, Ser.~I, Vol.~2.

\bibitem{weyl}
H.~Weyl, \emph{Symmetry} (Princeton, 1952).

\end{thebibliography}

\end{document}
